\chapter{Methods}

The probe-design method was implemented as a single, reproducible Python pipeline. Given a
target gene, the pipeline proceeds through the stages summarised below, each of which is
described in its own section together with the rationale for the choices made and the values of
the parameters used. All stages are driven by a small set of biologically motivated defaults
that can be overridden per organism type, per species and per gene (Section~\ref{sec:params}).

\begin{enumerate}
  \item retrieval and curation of gene sequences from \acrshort{ncbi};
  \item automatic length clustering;
  \item \acrfull{msa} and per-column conservation scoring;
  \item generation of sliding-window candidate probes;
  \item thermodynamic screening (\acrshort{tm}, GC, hairpin, homodimer);
  \item secondary-structure evaluation and No-fold scoring;
  \item organism-aware resolution of all thresholds;
  \item quality-based selection of the best candidates per gene;
  \item three-dimensional validation of the selected probes.
\end{enumerate}

\section{Sequence retrieval and curation}

For each target, gene sequences were retrieved programmatically from the \acrshort{ncbi}
nucleotide database through the Entrez interface (Biopython). Each target defines a primary
query combining the gene name, the organism and a sequence-length range, with progressively
broader alternative queries and a RefSeq accession as fallback, so that even poorly annotated
genes return a usable set. The number of sequences fetched per gene is configurable; a value
of $100$ was used in this work, justified by the rarefaction analysis of
Section~\ref{sec:diversity}.

Because public records mix full-length genes with partial fragments and oversized entries,
which together produce gap-ridden alignments, an automatic \emph{length-clustering} step
retains only the dominant length cluster: the densest sequence length is found and only
sequences within $\pm 25\%$ of it are kept. This selection is data-driven --- no per-gene
minimum or maximum length is tuned by hand --- so the method transfers to new genes without
manual adjustment.

\section{Alignment and conservation}

The curated sequences were aligned with MAFFT \citep{katoh2013} in its automatic mode; a
single-sequence mode is used when only one sequence is available. Conservation was then scored
per alignment column using the \acrfull{ppi} adopted from \emph{ViruScope} \citep{viruscope}:
the average identity over all pairs of sequences in a column, with weights that correctly
account for ambiguous IUPAC bases. The \acrshort{ppi} is a stricter and more informative
measure than the frequency of the most common base.

Unlike in primer design, the identity of the terminal $3'$ nucleotides (sometimes scored
separately as PPI3) was \emph{deliberately excluded}: capture probes are not extended by a
polymerase, so the fidelity of their $3'$ end is irrelevant to their function.

\section{Candidate generation}

Candidate probes were generated by sliding a window over the alignment. Windows of length
$18$--$28$ nucleotides \citep{wetmur1991} were taken every $3$ positions, and only windows whose
conservation exceeded the target's conservation threshold and whose gap fraction was below the
allowed maximum were kept for scoring.

\section{Thermodynamic screening}

Each candidate was screened with the nearest-neighbour thermodynamic model as implemented in
\texttt{primer3} \citep{untergasser2012,santalucia2004}. Four quantities were evaluated: the
melting temperature \acrshort{tm}, the GC content, the hairpin free energy and the homodimer
(self-dimer) free energy. A candidate passes the basic screen when all four fall within the
limits resolved for its target (Table~\ref{tab:params}). The \acrshort{tm} window can also be
derived automatically from the hybridisation assay temperature $T$, as $[T+15,\,T+35]\,^{\circ}$C,
so that the criterion tracks the real operating conditions of the device.

\begin{table}[H]
\begin{center}
\begin{tabular}{l l l}
\hline
\textbf{Parameter} & \textbf{Default value} & \textbf{Basis} \\
\hline
Probe length            & 18--28 nt              & \citet{wetmur1991} \\
Melting temperature $T_m$ & 53--72 $^{\circ}$C   & \citet{santalucia2004}; assay-derivable \\
GC content              & 40--60\%               & IDT / PremierBiosoft guidelines \\
Hairpin $\Delta G$      & $\geq -2.0$ kcal/mol   & IDT / primer3 guidance \\
Homodimer $\Delta G$    & $\geq -5.0$ kcal/mol   & IDT guidance \\
Conservation (\acrshort{ppi}) & $\geq 0.85$      & pan-strain detection \\
Alignment gap fraction  & $\leq 20\%$            & alignment quality \\
Secondary structure $\Delta G_{\text{MFE}}$ & $\geq -2.6$ kcal/mol & 5th percentile of the data \\
\hline
\end{tabular}
\end{center}
\caption{Global default thresholds. Values are overridden per organism type, species or gene
where justified (Section~\ref{sec:params}).}
\label{tab:params}
\end{table}

\section{Secondary structure and the No-fold score}

The single-strand secondary structure of each candidate was computed with \texttt{seqfold}
\citep{seqfold}, which returns the \acrfull{mfe} at $37\,^{\circ}$C and the fraction of paired
bases. Candidates whose \acrshort{mfe} was too negative (too stable a fold) were rejected. The
threshold of $-2.6$ kcal/mol corresponds to the fifth percentile of the observed
\acrshort{mfe} distribution; fixing it at this empirical value keeps the criterion reproducible
across runs, instead of always cutting a fixed fraction of the data.

In addition to the pass/fail screen, the three folding energies (hairpin, homodimer and
\acrshort{mfe}) were combined into a continuous \emph{No-fold} score in the range $0$--$100$
through a logistic transformation, again following \emph{ViruScope} \citep{viruscope}. A high
No-fold score denotes a probe that is largely free of stable self-structure. This continuous
score is used, alongside conservation, to rank the candidates rather than merely to filter
them.

\section{Organism-aware parameters}
\label{sec:params}

Because the intended application spans pathogens with very different genomes, a single global
set of thresholds is inadequate: the criteria appropriate for a balanced bacterial genome
would unfairly reject an AT-rich or a GC-rich target. All thresholds are therefore resolved
through a layered scheme, from the most specific to the most general:
\[
\text{gene} \;\rightarrow\; \text{species} \;\rightarrow\; \text{organism type}
\;\rightarrow\; \text{global default}.
\]
The first layer that defines a given parameter wins. Five organism types are distinguished,
differing mainly in genome GC content and in variability (Table~\ref{tab:types}); species-level
profiles then fine-tune the values using the GC content of each genome (for example
\emph{P. aeruginosa} at $\sim$67\% GC \citep{stover2000}, \emph{A. fumigatus} at $\sim$50\% GC
\citep{nierman2005}, and the AT-rich \emph{P. falciparum} at $\sim$19\% GC \citep{gardner2002}).

\begin{table}[H]
\begin{center}
\begin{tabular}{l c c c}
\hline
\textbf{Organism type} & \textbf{GC window} & \textbf{Conservation} & \textbf{Max length} \\
\hline
Bacteria    & 40--60\% & 0.85 & 28 \\
Virus       & 35--68\% & 0.70 & 30 \\
Fungus      & 40--60\% & 0.80 & 28 \\
Protozoan   & 15--40\% & 0.80 & 30 \\
Host (human)& 35--60\% & ---  & 28 \\
\hline
\end{tabular}
\end{center}
\caption{Per-type parameter profiles. Viruses use a wider GC window and a more permissive
conservation threshold to accommodate their variability; the protozoan profile is AT-shifted.}
\label{tab:types}
\end{table}

When the pipeline encounters an organism for which no profile exists, it does not silently fall
back to a default. If run interactively, it prompts the user in real time for the organism type
and for each threshold, proposing a suggested value (inferred from the organism name and its
genome composition) that can be accepted or overridden; the resulting profile is stored and
reused in later runs. When run non-interactively, it applies the suggested values without
blocking.

\section{Quality-based selection for 3D validation}

Three-dimensional prediction is comparatively expensive, so only the best candidates per gene
are carried forward. Rather than ranking by melting temperature --- which the thermodynamic
screen has already bounded --- the surviving candidates are ranked by a \emph{quality} score
that averages the two signals the pipeline trusts most, conservation and the No-fold score:
\[
\text{quality} \;=\; \tfrac{1}{2}\left(\text{\acrshort{ppi}} + \tfrac{\text{No-fold}}{100}\right).
\]
The top-ranked candidates of each gene are then exported for structural validation. The same
quality ranking is used to export larger curated sets (for example the top~25 per gene) for
downstream data-driven analysis.

\section{Three-dimensional validation}

The selected probes were submitted to Boltz-2 \citep{wohlwend2024}, run on a GPU, which
predicts the three-dimensional structure of each \acrshort{ssdna} sequence together with
confidence estimates: an overall confidence, the \acrfull{plddt} and the \acrfull{pae} matrix,
in the style of AlphaFold~3 \citep{abramson2024}. As these confidence metrics are calibrated
primarily for proteins, the overall confidence and the \acrshort{plddt} are used here as
relative, independent indicators of structural soundness of the short probes rather than as
absolute quality measures. The predictions were merged with the sequence metadata to produce a
final ranked shortlist combining sequence quality (conservation and No-fold) with the 3D
confidence.

\section{Integration of a reference probe panel}

An existing experimental probe panel (referred to as the reference panel) was integrated into
the same workflow for direct comparison. Each reference probe, supplied as a
(\texttt{job\_name}, \texttt{sequence}, \texttt{species}) record, was scored with exactly the
same metrics as the designed probes and by the criteria of its \emph{own} species and type.
Conservation is not defined for these probes, as they were not derived from an alignment and
therefore have no region of origin. The reference probes are stored alongside the designed
ones, with a source label, so that the two sets appear side by side with identical columns and
can be compared both in sequence space and, when submitted together, in 3D.

\section{Diversity and rarefaction analysis}
\label{sec:diversity}

To characterise the genetic diversity of each target and to justify the sampling depth, an
alignment-free analysis based on $k$-mer composition was performed. Pairwise cosine and Jaccard
distances between sequences quantify compositional and vocabulary diversity respectively,
while the fraction of unique sequences detects redundancy (exact duplicates, which are common
in \acrshort{ncbi}). A rarefaction analysis --- subsampling $N$ sequences per gene and measuring
$k$-mer diversity and richness --- was used to determine the value of $N$ beyond which little
new information is gained.

\section{Implementation and reproducibility}

The pipeline is implemented in Python, using Biopython for \acrshort{ncbi} access and sequence
handling, \texttt{primer3-py} for thermodynamics, \texttt{seqfold} for secondary structure, and
\texttt{pandas} for the tabular output; MAFFT is called as an external aligner and Boltz-2 is
run on a GPU. All intermediate and final results are written as versioned CSV files. Because
\acrshort{ncbi} does not always return an identical set of records, per-gene counts may vary by
a few percent between runs; for strict reproducibility the accession identifiers can be frozen.
